% !PRINTED 387
\weberpage{387}
%<<CONT
\weq{7}{
\begin{aligned}
&\left(\frac{a}{l}\right),\ \left(\frac{a}{l'}\right), \dots && m \equiv 3 \pmod 4, \\
&\left(\frac{-1}{a}\right),\ \left(\frac{a}{l}\right),\ \left(\frac{a}{l'}\right), \dots && m \equiv 1 \pmod 4, \\
& && m \equiv 4 \pmod 8, \\
&\left(\frac{2}{a}\right),\ \left(\frac{a}{l}\right),\ \left(\frac{a}{l'}\right), \dots && m \equiv 6 \pmod 8, \\
&\left(\frac{-2}{a}\right),\ \left(\frac{a}{l}\right),\ \left(\frac{a}{l'}\right), \dots && m \equiv 2 \pmod 8, \\
&\left(\frac{-1}{a}\right),\ \left(\frac{2}{a}\right),\ \left(\frac{a}{l}\right),\ \left(\frac{a}{l'}\right), \dots && m \equiv 0 \pmod 8.
\end{aligned}
}
Die Anzahl dieser Grundcharaktere ist, wie man sieht, gleich
der in \S\,104, 3. näher bestimmten Zahl $\lambda$ und die Anzahl der
Vorzeichenkombinationen ist gleich $2^\lambda$.

Die Abhängigkeit zwischen diesen Charakteren kann man so
darstellen.

Ist
\[
m = n^2 m', \quad D = Q^2\Delta = -4n^2m'
\]
und $n^2$ die größte in $m$ aufgehende Quadratzahl, so ist
\weq{8}{
\begin{aligned}
&\Delta \equiv -m', && \text{wenn } m' \equiv 3 \pmod 4, \\
&\Delta \equiv -4m' && \text{,,}\quad m' \equiv 1, 2 \pmod 4;
\end{aligned}
}
im ersten Fall ist $m'$ ungerade und nicht gleich 1, also gleich
einer der Primzahlen $l$ oder gleich einem Produkt aus mehreren
von ihnen, und aus
\[
\left(\frac{\Delta}{a}\right) = \left(\frac{-m'}{a}\right) = \left(\frac{a}{m'}\right) = 1,
\]
worin das Reziprozitätsgesetz der quadratischen Reste angewandt
ist, ergibt sich eine Relation zwischen den Charakteren.

Im zweiten Fall ist bei ungeradem $m'$
\[
\left(\frac{-m'}{a}\right) = \left(\frac{-1}{a}\right)\left(\frac{a}{m'}\right) = 1,
\]
und dies ist wieder ein Produkt mehrerer der Charaktere (7),
und ist endlich $m' = 2m''$ gerade, so ist
\[
\left(\frac{-m'}{a}\right) = \left(\frac{\pm 2}{a}\right)\left(\frac{a}{m''}\right),
\]
wo das Zeichen $\pm$ so bestimmt wird, daß $\pm m'' \equiv 1 \pmod 4$ wird.
Also ist wiederum die Anzahl der Geschlechter höchstens $= 2^{\lambda-1}$.
